%! Unit 2 — Exponential and Logarithmic Functions — Unit Cover Page
\documentclass[10pt]{article}
\usepackage{precalculus-article}
\usepackage{precalculus-boxes}
\usepackage{ltablex}
\keepXColumns

\begin{document}
\thispagestyle{empty}

% Full-bleed navy banner
\begin{tikzpicture}[remember picture, overlay]
  \fill[navy] (current page.north west)
    rectangle ([yshift=-1.15in]current page.north east);
\end{tikzpicture}
\vspace{-1.05in}

\begin{center}
  {\color{white}\textbf{\LARGE Precalculus \quad 2026--2027}}\\[6pt]
  {\color{skymid}\Large\textbf{Unit 2: Exponential and Logarithmic Functions}}\\[4pt]
  {\color{white}\small AP Exam Weighting: 27--40\% \quad|\quad Class Periods: 30--45}
\end{center}

\vspace{0.26in}

\begin{tcolorbox}[colback=goldbg, colframe=goldacc, boxrule=0.8pt, arc=2mm,
    left=3mm, right=3mm, top=2mm, bottom=2mm]
\small\textbf{Unit Essential Questions:}
\begin{itemize}[leftmargin=1.4em, itemsep=2pt, topsep=2pt]
  \item How do exponential and logarithmic functions model situations of growth, decay, and change over time?
  \item What is the relationship between exponential and logarithmic functions, and how does that relationship inform algebraic and graphical reasoning?
\end{itemize}
\end{tcolorbox}

\vspace{0.14in}

\begin{tocbox}
\small
\renewcommand{\arraystretch}{1.45}
\begin{tabularx}{\linewidth}{@{} c >{\bfseries}p{0.52\linewidth} l X @{}}
\rowcolor{sky}
\textbf{\#} & \textbf{Lesson Title} & \textbf{AP Topic} & \textbf{AP Skill} \\
\midrule
2.1  & Change in Arithmetic and Geometric Sequences                & 2.1  & 1.C, 3.A \\
2.2  & Change in Linear and Exponential Functions                  & 2.2  & 1.A, 3.A \\
2.3  & Exponential Functions                                       & 2.3  & 3.C \\
2.4  & Exponential Function Manipulation                           & 2.4  & 1.A, 3.A \\
2.5  & Exponential Function Context and Data Modeling              & 2.5  & 1.C, 3.A \\
2.6  & Competing Function Model Validation                         & 2.6  & 2.A, 3.C \\
2.7  & Composition of Functions                                    & 2.7  & 1.A, 2.B \\
2.8  & Inverse Functions                                           & 2.8  & 1.B, 2.B \\
2.9  & Logarithmic Expressions                                     & 2.9  & 1.A \\
2.10 & Inverses of Exponential Functions                           & 2.10 & 1.B, 2.B \\
2.11 & Logarithmic Functions                                       & 2.11 & 3.C \\
2.12 & Logarithmic Function Manipulation                           & 2.12 & 1.A, 1.B \\
2.13 & Exponential and Logarithmic Equations and Inequalities      & 2.13 & 1.A \\
2.14 & Logarithmic Function Context and Data Modeling              & 2.14 & 1.C, 3.A \\
2.15 & Semi-log Plots                                              & 2.15 & 2.B, 3.C \\
\midrule
\rowcolor{sky}
     & \textbf{Sample Test}                                        &      &       \\
\end{tabularx}
\end{tocbox}

\vspace{0.14in}

\begin{skillbox}[AP Mathematical Practices --- Unit 2 Focus]{goldbox}
\renewcommand{\arraystretch}{1.4}
\begin{tabularx}{\linewidth}{@{} >{\bfseries\color{navy}}p{0.12\linewidth} X @{}}
MP 1.A & Solve equations and inequalities analytically, with and without technology. \\
MP 1.B & Express functions, equations, or expressions in analytically equivalent forms. \\
MP 1.C & Construct new functions using transformations, compositions, inverses, or regressions. \\
MP 2.A & Identify information from graphical, numerical, analytical, and verbal representations. \\
MP 2.B & Construct equivalent graphical, numerical, analytical, and verbal representations. \\
MP 3.A & Describe the characteristics of a function with varying levels of precision. \\
MP 3.C & Support conclusions or choices with a logical rationale or appropriate data. \\
\end{tabularx}
\end{skillbox}

\end{document}
